\scanpage{28}
dans $\mathcal E$, qu'on peut choisir canoniquement en prenant le quart de grand cercle qui joint $P$ au milieu $2$ du \deleted{segment} l'arc $(1,\infty)$ de grand cercle opposé à $0$; ce qui définit un \unsure{isom} du rev. universel de $U$ en $P$, avec celui de $U$ en $\bar P$ sur lequel le groupe précédent opère de façon canonique, donc aussi par transport de structure sur $\tilde U\deleted{\unsure{_P}}(P)$, donc il s'envoie dans $\mathcal E$.

On trouve que
\[
\dot\sigma_0\mapsto\sigma_0,\quad\dot\tau\to\tau_0
\]
est le scindage correspondant. De $=$ pour les pts fixes $-1$ de $\dot\sigma_1$, ou $\frac12$ de $\dot\sigma_\infty$, on trouve des scindages partiels sur les ss-groupes de \deleted{\unsure{conj.}} stabilité $\Gamma_{-1}\subset\Gamma$ et $\Gamma_2\subset\Gamma$. Il est commode de poser
\[
\deleted{\unsure{\Gamma}}\quad 2=0',\quad -1=1',\quad \tfrac12=\infty',
\]
d'où des groupes d'isotropie $\Gamma_{0'}=\{\dot\sigma_0,\tau\}$, $\Gamma_{1'}=\{\dot\sigma_1,\tau\}$, et $\Gamma_{\infty'}=\{\dot\sigma_\infty,\tau\}$, qui se \unsure{remontent} resp. en $\sigma_0$ et $\tau_0$, \unsure{ou} $\sigma_1$ et $\tau_1$, ou $\sigma_\infty$ et $\tau_\infty$.

Les composantes connexes de $U^\tau$ sont les trois arcs de grand cercle entre $0,1,\infty$, qui contiennent chacun un \deleted{pt} et un seul des pts $0',1',\infty'$, et les classes de conjugaison
\mspage{27}
